\chapter{Практика 4. QPSK модуляция}

\section{Задание}
Реализовать операцию QPSK модуляции битового сообщения и QPSK демодуляции символов.

\section{Теоретическое описание}
QPSK (Quadrature Phase Shift Keying) — метод фазовой модуляции, при котором фаза несущего сигнала принимает одно из четырех возможных значений. Каждому значению фазы соответствует пара бит (дибит). Скорость передачи символов в два раза ниже скорости передачи бит.

Таблица QPSK модуляции:

\begin{center}
\begin{tabular}{|c|c|c|c|}
\hline
\textbf{Дибит} & \textbf{I} & \textbf{Q} & \textbf{Символ} \\
\hline
00 & +0.707 & +0.707 & $0.707 + 0.707j$ \\
01 & +0.707 & -0.707 & $0.707 - 0.707j$ \\
10 & -0.707 & +0.707 & $-0.707 + 0.707j$ \\
11 & -0.707 & -0.707 & $-0.707 - 0.707j$ \\
\hline
\end{tabular}
\end{center}

\section{Реализация}

\begin{lstlisting}[language=Python, caption=QPSK модулятор и демодулятор]
class Modulator:
    def modulate(bits):
        if not bits:
            return []
        if len(bits) % 2 != 0:
            bits = bits + '0'
        symbols = []
        for i in range(0, len(bits), 2):
            bit_pair = bits[i:i + 2]
            if bit_pair == '00':
                symbols.append(complex(0.707, 0.707))
            elif bit_pair == '01':
                symbols.append(complex(0.707, -0.707))
            elif bit_pair == '10':
                symbols.append(complex(-0.707, 0.707))
            else:
                symbols.append(complex(-0.707, -0.707))
        return symbols

class Demodulator:
    def demodulate(symbols):
        if not symbols:
            return ""
        bits = []
        for symbol in symbols:
            if symbol.real > 0 and symbol.imag > 0:
                bits.append('00')
            elif symbol.real > 0 and symbol.imag < 0:
                bits.append('01')
            elif symbol.real < 0 and symbol.imag > 0:
                bits.append('10')
            else:
                bits.append('11')
        return ''.join(bits)
\end{lstlisting}

\section{Результат}
\begin{verbatim}
QPSK модуляция:
  До модуляции: 570 бит
  После модуляции: 285 символов
\end{verbatim}
Демодуляция выполняется путем определения области комплексной плоскости, в которую попал принятый символ.

\section{Вывод}
QPSK модулятор и демодулятор успешно реализованы. При отсутствии шума демодуляция выполняется без ошибок.